\documentclass[11pt,a4paper]{amsart}

\usepackage[margin=1in]{geometry}
\usepackage[T1]{fontenc}
\usepackage{lmodern}
\usepackage{microtype}
\usepackage{amsmath,amssymb,amsthm,mathtools}
\usepackage{enumitem}
\usepackage{xcolor}
\usepackage[colorlinks=true,linkcolor=blue!60!black,
  citecolor=blue!60!black,urlcolor=blue!60!black]{hyperref}

\allowdisplaybreaks

% ─── theorem environments ───
\theoremstyle{plain}
\newtheorem{theorem}{Theorem}[section]
\newtheorem{lemma}[theorem]{Lemma}
\newtheorem{proposition}[theorem]{Proposition}
\newtheorem{corollary}[theorem]{Corollary}

\theoremstyle{definition}
\newtheorem{definition}[theorem]{Definition}

\theoremstyle{remark}
\newtheorem{remark}[theorem]{Remark}

% ─── macros ───
\newcommand{\R}{\mathbb{R}}
\newcommand{\PnR}[1][n]{\mathcal{P}^{\R}_{#1}}
\newcommand{\fp}{\boxplus_n}
\DeclareMathOperator{\tr}{tr}
\DeclareMathOperator{\disc}{disc}
\DeclareMathOperator{\rank}{rank}

\begin{document}

\title{The Finite Free Stam Inequality}

\date{\today}

\begin{abstract}
The classical Stam inequality asserts the superadditivity of
the reciprocal Fisher information under convolution of
independent random variables.
We prove the polynomial analogue in the framework of finite
free probability: for monic, degree-$n$, real-rooted
polynomials~$p$ and~$q$,
\[
  \frac{1}{\Phi_n(p\fp q)}
  \;\ge\;
  \frac{1}{\Phi_n(p)}+\frac{1}{\Phi_n(q)},
\]
where $\fp$ is the symmetric additive convolution of Marcus,
Spielman, and Srivastava, and $\Phi_n$ is the finite free
Fisher information.
The proof combines an algebraic inequality---the
Score-Gradient Inequality, established via two applications
of Cauchy--Schwarz---with a flow-based argument exploiting
the Hermite semigroup structure of~$\fp$, together with
the Jacobian contraction inequality for the root map.
We also derive a closed-form expression for $\Phi_n$ via
residue calculus and verify the inequality explicitly for
cubics.
\end{abstract}

\maketitle
\tableofcontents

% ═══════════════════════════════════════════════════════════
\section{Introduction}
% ═══════════════════════════════════════════════════════════

\subsection{Background and motivation}

In information theory, the Stam inequality~\cite{stam}
states that if $X$ and $Y$ are independent random variables
with finite Fisher information $I(X)$ and $I(Y)$, then
\[
  \frac{1}{I(X+Y)}
  \;\ge\;\frac{1}{I(X)}+\frac{1}{I(Y)}.
\]
This fundamental inequality---equivalent to the entropy power
inequality of Shannon and Stam---captures the principle that
convolution of independent sources strictly increases disorder.

Finite free probability, introduced by Marcus, Spielman,
and Srivastava~\cite{MSS}, provides a polynomial analogue
of free probability in which random variables are replaced
by real-rooted polynomials and addition by a
deterministic convolution operation~$\fp$.
Within this framework, the natural question arises:

\medskip
\emph{Does the Stam inequality hold for the finite free
additive convolution?}
\medskip

The purpose of this paper is to answer this question
affirmatively.

\subsection{Statement of the main result}

Let $\mathcal{P}_n$ denote the space of monic polynomials of
degree~$n$ with real coefficients, and
$\PnR\subset\mathcal{P}_n$ the subset with all real roots.
For $p\in\PnR$ with distinct roots
$\lambda_1<\cdots<\lambda_n$, define the \emph{scores}
\[
  V_i:=\sum_{j\ne i}\frac{1}{\lambda_i-\lambda_j}
\]
and the \emph{finite free Fisher information}
$\Phi_n(p):=\sum_{i=1}^n V_i^2$.
The symmetric additive convolution $p\fp q$ is recalled in
Section~2.

\begin{theorem}[Finite Free Stam Inequality]
\label{thm:main}
For $p,q\in\PnR$ with positive variance,
\begin{equation}\label{eq:stam}
  \frac{1}{\Phi_n(p\fp q)}
  \;\ge\;\frac{1}{\Phi_n(p)}+\frac{1}{\Phi_n(q)}.
\end{equation}
\end{theorem}

The proof combines four ingredients: the Score-Gradient
Inequality (Theorem~\ref{thm:sgi}), a dissipation identity
for the Hermite flow (Lemma~\ref{lem:dissipation}),
the Score--Jacobian relation
(Proposition~\ref{prop:score-jac}), and the Jacobian
contraction inequality
(Proposition~\ref{prop:jac-contract}).

\medskip\noindent
\textbf{Convention.}
We set $\Phi_n(p):=\infty$ (equivalently
$1/\Phi_n(p):=0$) when $p$ has a repeated root.
The proof below establishes~\eqref{eq:stam} under
the hypothesis that $p$ and~$q$ have distinct roots;
the extension to the boundary of $\PnR$ is discussed
in Remark~\ref{rem:boundary}.


% ═══════════════════════════════════════════════════════════
\section{Preliminaries}
% ═══════════════════════════════════════════════════════════

\subsection{Root statistics}

For $p(x)=\prod_{i=1}^n(x-\lambda_i)
=\sum_{k=0}^n a_k x^{n-k}$ with $a_0=1$, the
\emph{mean} and \emph{variance} of the root distribution are
\[
  \bar\lambda:=\frac{1}{n}\sum_{i=1}^n\lambda_i,
  \qquad
  \sigma^2(p):=\frac{1}{n}\sum_{i=1}^n
    (\lambda_i-\bar\lambda)^2.
\]

\begin{lemma}\label{lem:var-coeff}
$\sigma^2(p)=\dfrac{(n-1)a_1^2}{n^2}-\dfrac{2a_2}{n}$.
\end{lemma}

\begin{proof}
By Vieta's formulas, $\sum_i\lambda_i=-a_1$ and
$\sum_{i<j}\lambda_i\lambda_j=a_2$, whence
$\sum_i\lambda_i^2=a_1^2-2a_2$.
The result follows from
$\sigma^2=\frac{1}{n}\sum_i\lambda_i^2-\bar\lambda^2$.
\end{proof}

\subsection{Symmetric additive convolution}

The \emph{finite free additive convolution} is defined by
averaging over the orthogonal group: if $A$ and $B$ are real
symmetric matrices with characteristic polynomials~$p$
and~$q$, then
\[
  (p\fp q)(x)
  =\int_{O(n)}\det\bigl(xI-(A+QBQ^T)\bigr)\,d\mu(Q).
\]
By the MSS theorem~\cite{MSS}, this admits a differential
operator representation: writing
$q(x)=\sum_{k=0}^n b_k x^{n-k}$,
\begin{equation}\label{eq:diff-op}
  (p\fp q)(x)=T_q\,p(x),
  \qquad
  T_q:=\sum_{k=0}^n\frac{(n-k)!}{n!}\,b_k\,\partial_x^k.
\end{equation}
The coefficients of $r=p\fp q$,
$r(x)=\sum_k c_k x^{n-k}$, satisfy
\begin{equation}\label{eq:ck}
  c_k=\sum_{i+j=k}
  \frac{(n-i)!\,(n-j)!}{n!\,(n-k)!}\,a_i\,b_j.
\end{equation}

Two fundamental properties:

\begin{theorem}[\cite{MSS}]\label{thm:mss}
If $p,q\in\PnR$, then $p\fp q\in\PnR$.
\end{theorem}

\begin{lemma}[Variance additivity]\label{lem:var-add}
$\sigma^2(p\fp q)=\sigma^2(p)+\sigma^2(q)$.
\end{lemma}

\begin{proof}
From~\eqref{eq:ck},
$c_1=a_1+b_1$ and
$c_2=a_2+\frac{n-1}{n}a_1 b_1+b_2$.
Substituting into Lemma~\ref{lem:var-coeff} and expanding
$(a_1+b_1)^2$, the cross-terms
$\frac{2(n-1)a_1 b_1}{n^2}$ and
$-\frac{2(n-1)a_1 b_1}{n^2}$ cancel.
\end{proof}


\subsection{Scores and Fisher information}

\begin{definition}\label{def:scores}
For $p\in\PnR$ with distinct roots
$\lambda_1<\cdots<\lambda_n$, the \emph{score} at~$\lambda_i$
is $V_i:=\sum_{j\ne i}(\lambda_i-\lambda_j)^{-1}$, the
\emph{Fisher information} is $\Phi_n(p):=\sum_i V_i^2$, and
the \emph{score-gradient energy} is
$\mathcal{S}(p):=\sum_{i<j}(V_i-V_j)^2/(\lambda_i-\lambda_j)^2$.
\end{definition}

\begin{lemma}\label{lem:score-Vi}
$V_i=p''(\lambda_i)/(2\,p'(\lambda_i))$.
\end{lemma}

\begin{proof}
Since $p'(\lambda_i)=\prod_{j\ne i}(\lambda_i-\lambda_j)$,
differentiating once more gives
$p''(\lambda_i)=2\sum_{k\ne i}\prod_{j\ne i,\,j\ne k}
(\lambda_i-\lambda_j)=2\,p'(\lambda_i)\,V_i$.
\end{proof}

\begin{lemma}[Score identities]\label{lem:score-ids}
\begin{enumerate}[label=\textup{(\roman*)},nosep]
\item $\sum_i V_i=0$.
\item $\sum_i\lambda_i V_i=\binom{n}{2}$.
\item $\sum_i(\lambda_i-\bar\lambda)V_i=\binom{n}{2}$.
\item $\Phi_n(p)
  =\sum_{i<j}\dfrac{V_i-V_j}{\lambda_i-\lambda_j}$.
\end{enumerate}
\end{lemma}

\begin{proof}
(i): $\sum_i V_i=\sum_{i\ne j}(\lambda_i-\lambda_j)^{-1}=0$
by antisymmetry.

(ii): $\sum_i\lambda_i V_i
=\sum_{i\ne j}\lambda_i/(\lambda_i-\lambda_j)
=\sum_{i<j}\bigl[\lambda_i/(\lambda_i-\lambda_j)
  +\lambda_j/(\lambda_j-\lambda_i)\bigr]
=\sum_{i<j}1=\binom{n}{2}$.

(iii): Immediate from (ii) and (i).

(iv): $\sum_i V_i^2
=\sum_{i\ne j}V_i/(\lambda_i-\lambda_j)
=\sum_{i<j}(V_i-V_j)/(\lambda_i-\lambda_j)$.
\end{proof}

\begin{lemma}[Fisher--variance inequality]
\label{lem:fisher-var}
$\Phi_n(p)\cdot\sigma^2(p)\ge n(n-1)^2/4$.
\end{lemma}

\begin{proof}
By Cauchy--Schwarz applied to
Lemma~\ref{lem:score-ids}(iii):
$\binom{n}{2}^2
\le\bigl(\sum(\lambda_i-\bar\lambda)^2\bigr)
  \bigl(\sum V_i^2\bigr)
=n\,\sigma^2\cdot\Phi_n$.
\end{proof}


% ═══════════════════════════════════════════════════════════
\section{The Score-Gradient Inequality}
\label{sec:sgi}
% ═══════════════════════════════════════════════════════════

\begin{theorem}[Score-Gradient Inequality]
\label{thm:sgi}
For $p\in\PnR$ of degree $n\ge 2$ with distinct roots,
\begin{equation}\label{eq:sgi}
  \mathcal{S}(p)\,\sigma^2(p)
  \;\ge\;\frac{n-1}{2}\,\Phi_n(p),
\end{equation}
with equality if and only if $V_i=c(\lambda_i-\bar\lambda)$
for some constant~$c$.
\end{theorem}

\begin{proof}
Write $T=n\,\sigma^2$, $U=\Phi_n$, $S=\mathcal{S}$.
The claim is $S\,T\ge\frac{n(n-1)}{2}\,U$.

\emph{Step~1.}
By Lemma~\ref{lem:score-ids}(iii) and Cauchy--Schwarz,
\begin{equation}\label{eq:cs1}
  \frac{n^2(n-1)^2}{4}\;\le\;T\,U.
\end{equation}

\emph{Step~2.}
By Lemma~\ref{lem:score-ids}(iv) and Cauchy--Schwarz,
\begin{equation}\label{eq:cs2}
  U^2\;\le\;S\cdot\binom{n}{2}.
\end{equation}

\emph{Step~3.}
Combining:
$S\,T
\ge\frac{2\,U^2}{n(n-1)}\cdot T
=\frac{2\,U}{n(n-1)}\cdot T\,U
\ge\frac{2\,U}{n(n-1)}\cdot\frac{n^2(n-1)^2}{4}
=\frac{n(n-1)}{2}\,U$.

\emph{Equality.}
Equality in~\eqref{eq:cs1} holds iff
$V_i=c(\lambda_i-\bar\lambda)$.
Equality in~\eqref{eq:cs2} holds iff
$(V_i-V_j)/(\lambda_i-\lambda_j)$ is constant.
These conditions are equivalent: if
$V_i=c(\lambda_i-\bar\lambda)$, then
$(V_i-V_j)/(\lambda_i-\lambda_j)=c$ for all $i<j$.
\end{proof}

\begin{remark}\label{rem:hermite}
The equality condition $V_i=c(\lambda_i-\bar\lambda)$
characterises, up to affine transformation, the zeros of
the Hermite polynomial~$H_n$: evaluating the ODE
$H_n''-2xH_n'+2nH_n=0$ at a zero~$x_k$ yields $V_k=x_k$.
For $n=2$ this holds for all distinct root configurations.
\end{remark}


% ═══════════════════════════════════════════════════════════
\section{The Hermite flow}
\label{sec:flow}
% ═══════════════════════════════════════════════════════════

\subsection{The Gaussian semigroup}

\begin{definition}\label{def:gaussian}
For $t>0$, the \emph{finite Gaussian} $g_t\in\PnR$ is the
polynomial whose roots are at positions
$\sqrt{t}\,\xi_1<\cdots<\sqrt{t}\,\xi_n$, where
$\xi_1<\cdots<\xi_n$ are the zeros of the probabilist
Hermite polynomial~$\mathrm{He}_n$.
Thus $\sigma^2(g_t)=t$.
The semigroup property $g_s\fp g_t=g_{s+t}$ holds, and
$g_0:=x^n$.
\end{definition}

Fix $p,q\in\PnR$ with $a:=\sigma^2(p)>0$ and
$b:=\sigma^2(q)>0$.

\begin{definition}[Hermite flow]\label{def:flow}
Define $p_t:=p\fp g_{bt}$ for $t\ge 0$.
Then:
\begin{enumerate}[label=\textup{(\alph*)},nosep]
\item $p_0=p$;
\item $\sigma^2(p_t)=a+bt$ (by variance additivity,
  Lemma~\ref{lem:var-add});
\item $p_t\in\PnR$ for all $t\ge 0$
  (by Theorem~\ref{thm:mss}, since $g_{bt}\in\PnR$).
\end{enumerate}
\end{definition}

\begin{remark}\label{rem:why-gaussian}
In an earlier version of this paper, the flow was defined
as $p_t=p\fp q_t$ where $q_t$ is the fractional power
of~$q$ in the $K$-transform ($K_{q_t}=K_q^t$);
this gives $q_0=x^n$, $q_1=q$, and $\sigma^2(q_t)=tb$.
However, $q_t$ need not lie in~$\PnR$ for non-integer~$t$,
so $p_t$ is not guaranteed to be real-rooted along the flow,
and the root ODE~\eqref{eq:root-ode} can fail.
The Gaussian semigroup avoids this: $g_{bt}\in\PnR$
for all~$t$, the root ODE holds exactly, and
all dissipation identities follow.
The cost is that the flow reaches $p\fp g_b$ at $t=1$,
not $p\fp q$.
To bridge this gap, we introduce the
Score--Jacobian relation and Jacobian contraction
in Section~\ref{sec:jacobian}.
\end{remark}


\subsection{Perturbation analysis}

Since $g_{bt}\in\PnR$ and $p_t\in\PnR$ has simple roots for
$t=0$, a Lyapunov argument establishes simple roots for all
$t\ge 0$; we state the key steps.

\begin{lemma}\label{lem:heat}
The Hermite flow satisfies the backwards heat equation
\[
  \partial_t p_t(x)
  =-\frac{b}{2(n-1)}\,p_t''(x).
\]
\end{lemma}

\begin{proof}
The operator $T_{g_{bt}}$ acts via~\eqref{eq:diff-op}
with coefficients of~$g_{bt}$.
By the semigroup property $p_{t+h}=p_t\fp g_{bh}$, and
the $K$-transform of $g_{bh}$ equals
$\exp(-bh\,z^2/(2(n-1)))\bmod z^{n+1}$,
so the coefficient ODE
$\kappa_k'=-b/(2(n-1))\cdot\kappa_{k-2}$
translates to the stated PDE.
\end{proof}

\begin{lemma}[Root ODE]\label{lem:root-ode}
Let $\lambda_i(t)$ denote the roots of~$p_t$.
Then
\begin{equation}\label{eq:root-ode}
  \dot\lambda_i(t)
  =\frac{b}{n-1}\,V_i(t)
  =\frac{b}{n-1}\sum_{j\ne i}
   \frac{1}{\lambda_i(t)-\lambda_j(t)}.
\end{equation}
\end{lemma}

\begin{proof}
Differentiating $p_t(\lambda_i(t))=0$ in~$t$:
$\partial_t p_t(\lambda_i)+p_t'(\lambda_i)\dot\lambda_i=0$.
By Lemma~\ref{lem:heat},
$\partial_t p_t(\lambda_i)
=-\frac{b}{2(n-1)}\,p_t''(\lambda_i)$,
so $\dot\lambda_i
=\frac{b}{2(n-1)}\,
 \frac{p_t''(\lambda_i)}{p_t'(\lambda_i)}
=\frac{b}{n-1}\,V_i$
by Lemma~\ref{lem:score-Vi}.
\end{proof}

\begin{lemma}[Simple roots]\label{lem:simple}
For all $t\ge 0$, $p_t$ has $n$ simple real roots.
\end{lemma}

\begin{proof}
Real-rootedness holds by Theorem~\ref{thm:mss}.
For simplicity, define the log-Vandermonde
$W(t):=\sum_{i<j}\log(\lambda_j(t)-\lambda_i(t))$.
By Lemma~\ref{lem:root-ode},
\[
  \dot W(t)=\sum_{i<j}
  \frac{\dot\lambda_j-\dot\lambda_i}
       {\lambda_j-\lambda_i}
  =\frac{b}{n-1}\sum_{i<j}
  \frac{V_j-V_i}{\lambda_j-\lambda_i}
  =\frac{b}{n-1}\,\Phi_n(p_t)\ge 0,
\]
using Lemma~\ref{lem:score-ids}(iv).
Hence $W(t)\ge W(0)$, i.e.\
$\prod_{i<j}(\lambda_j(t)-\lambda_i(t))
\ge\prod_{i<j}(\lambda_j(0)-\lambda_i(0))=:D_0>0$.
Since $\sigma^2(p_t)=a+bt$ is bounded on compact
intervals, each pairwise gap has a uniform positive
lower bound, so the roots remain simple.
\end{proof}


\subsection{Dissipation and the integral identity}

\begin{lemma}[Dissipation]\label{lem:dissipation}
$\displaystyle\frac{d}{dt}\Phi_n(p_t)
=-\frac{2b}{n-1}\,\mathcal{S}(p_t)$.
\end{lemma}

\begin{proof}
From Lemma~\ref{lem:root-ode}, write
$\varepsilon:=bh/(n-1)$ for a small increment~$h$.
The perturbed scores satisfy
\[
  V_i^{(h)}
  =V_i-\varepsilon\sum_{j\ne i}
   \frac{V_i-V_j}{(\lambda_i-\lambda_j)^2}+O(h^2).
\]
Squaring and summing, then pairing $(i,j)$ with $(j,i)$:
\[
  \Phi_n(p_{t+h})=\Phi_n(p_t)
  -2\varepsilon\sum_{i<j}
   \frac{(V_i-V_j)^2}{(\lambda_i-\lambda_j)^2}+O(h^2)
  =\Phi_n(p_t)-\frac{2bh}{n-1}\,\mathcal{S}(p_t)+O(h^2).
\]
Taking $h\to 0$ yields the result
(the error terms are uniform by Lemma~\ref{lem:simple}).
\end{proof}

\begin{corollary}[Integral identity]\label{cor:integral}
\[
  \frac{1}{\Phi_n(p\fp g_b)}
  -\frac{1}{\Phi_n(p)}
  =\frac{2b}{n-1}\int_0^1
   \frac{\mathcal{S}(p_t)}{\Phi_n(p_t)^2}\,dt.
\]
\end{corollary}

\begin{proof}
Set $f(t):=1/\Phi_n(p_t)$.
By the chain rule and Lemma~\ref{lem:dissipation},
$f'(t)=\frac{2b}{n-1}\cdot
 \frac{\mathcal{S}(p_t)}{\Phi_n(p_t)^2}\ge 0$.
Now integrate from $0$ to~$1$.
\end{proof}


\subsection{Gaussian-input Stam inequality}

\begin{theorem}[Gaussian-input Stam]
\label{thm:gauss-stam}
For all $p\in\PnR$ with $\sigma^2(p)=a>0$ and $t>0$:
\begin{equation}\label{eq:gauss-stam}
  \frac{1}{\Phi_n(p\fp g_t)}
  \;\ge\;\frac{1}{\Phi_n(p)}+\frac{1}{\Phi_n(g_t)}.
\end{equation}
\end{theorem}

\begin{proof}
\emph{Step~1 (Differential inequality).}
The SGI (Theorem~\ref{thm:sgi}) applied to $p_s=p\fp g_{ts}$
gives $\mathcal{S}(p_s)\ge(n-1)\Phi_n(p_s)/(2\sigma^2(p_s))$.
Substituting into Lemma~\ref{lem:dissipation}
(with the Gaussian variance parameter~$t$):
\[
  \frac{d}{ds}\Phi_n(p_s)
  \;\le\;-\frac{t}{a+ts}\,\Phi_n(p_s).
\]
Integrating $(\log\Phi_n)'\le-t/(a+ts)$ from $0$ to~$s$:
\begin{equation}\label{eq:upper-bound}
  \frac{1}{\Phi_n(p_s)}
  \;\ge\;\frac{a+ts}{a\,\Phi_n(p)}.
\end{equation}

\emph{Step~2 (Integration).}
Substituting~\eqref{eq:upper-bound} into
Corollary~\ref{cor:integral} with $b=t$:
\[
  \frac{1}{\Phi_n(p\fp g_t)}-\frac{1}{\Phi_n(p)}
  \;\ge\;
  \int_0^1\frac{t}{(a+ts)\,\Phi_n(p_s)}\,ds
  \;\ge\;\int_0^1\frac{t}{a\,\Phi_n(p)}\,ds
  \;=\;\frac{t}{a\,\Phi_n(p)}.
\]
Wait---this gives
$1/\Phi_n(p\fp g_t)\ge(a+t)/(a\,\Phi_n(p))$, which is just~\eqref{eq:upper-bound} at $s=1$.
To get~\eqref{eq:gauss-stam}, we use the
lower bound on the dissipation rate directly.

From Corollary~\ref{cor:integral} and the SGI:
\[
  f'(s)
  =\frac{2t}{n-1}\cdot
   \frac{\mathcal{S}(p_s)}{\Phi_n(p_s)^2}
  \;\ge\;\frac{t}{\sigma^2(p_s)\,\Phi_n(p_s)}
  \;\ge\;\frac{t}{\sigma^2(p_s)}\cdot
         \frac{a+ts}{a\,\Phi_n(p)}.
\]
But we can also bound $f'(s)$ from below without
using~\eqref{eq:upper-bound}: the Cauchy--Schwarz inequality
in the $L$-inner product gives
$\mathcal{S}/\Phi_n^2\ge 1/\binom{n}{2}$
(as in the SGI proof, Step~2), hence
\[
  f'(s)\;\ge\;\frac{2t}{n-1}\cdot
  \frac{1}{\binom{n}{2}}
  =\frac{4t}{n(n-1)^2}.
\]
Integrating from $0$ to $1$:
\[
  \frac{1}{\Phi_n(p\fp g_t)}-\frac{1}{\Phi_n(p)}
  \;\ge\;\frac{4t}{n(n-1)^2}
  =\frac{1}{\Phi_n(g_t)},
\]
since $\Phi_n(g_t)=n(n-1)^2/(4t)$ (the scores of the
finite Gaussian satisfy $V_i=\lambda_i/(2t^{1/2}\cdot
t^{-1/2})=\lambda_i/2$\ldots more precisely,
by the Hermite ODE $\mathrm{He}_n''(x_k)=x_k\,
\mathrm{He}_n'(x_k)$ and scaling,
$\Phi_n(g_t)\cdot\sigma^2(g_t)=n(n-1)^2/4$, i.e.\
equality in the Fisher--variance inequality).
\end{proof}


% ═══════════════════════════════════════════════════════════
\section{The root map and Jacobian contraction}
\label{sec:jacobian}
% ═══════════════════════════════════════════════════════════

The Hermite flow proves the Stam inequality when one input
is Gaussian.
To handle general~$q$, we need a non-flow ingredient:
the \emph{Jacobian contraction} of the root map, which
relates the score vectors of $p$, $q$, and $p\fp q$
directly.

\subsection{The root map}

\begin{definition}\label{def:root-map}
For $p,q\in\PnR$ with ordered root vectors
$\alpha=(\alpha_1\ge\cdots\ge\alpha_n)$ and
$\beta=(\beta_1\ge\cdots\ge\beta_n)$, define
$\Omega(\alpha,\beta)\in\R^n$ to be the ordered root
vector of $p\fp q$.
Write $\gamma:=\Omega(\alpha,\beta)$.
Let $J$ denote the $n\times 2n$ Jacobian matrix of
$\Omega$ at $(\alpha,\beta)$, defined whenever the roots
are simple.
\end{definition}

\subsection{Score--Jacobian relation}

\begin{proposition}[Score--Jacobian relation]
\label{prop:score-jac}
Let $p,q\in\PnR$ have simple roots with ordered root
vectors $\alpha,\beta$ and
$\gamma=\Omega(\alpha,\beta)$.
For any $a,b\ge 0$:
\begin{equation}\label{eq:score-jac}
  J\begin{pmatrix}a\,V(\alpha)\\b\,V(\beta)\end{pmatrix}
  =(a+b)\,V(\gamma),
\end{equation}
where $V(\alpha)
=\bigl(\sum_{j\ne i}(\alpha_i-\alpha_j)^{-1}\bigr)_{i=1}^n$.
\end{proposition}

\begin{proof}
For $t\ge 0$, define
$p_t:=\exp(-\tfrac{t}{2}\partial_x^2)\,p$
with ordered root vector $\alpha(t)$, and similarly
$q_t,\,r_t$ with $\beta(t),\,\gamma(t)$.
Since $\fp$ commutes with any differential operator
(because the $K$-transform converts differentiation
to multiplication),
$r_{(a+b)t}=p_{at}\fp q_{bt}$, whence
$\gamma((a+b)t)=\Omega(\alpha(at),\beta(bt))$.
Differentiating at $t=0$:
\[
  (a+b)\,\gamma'(0)
  =J\begin{pmatrix}a\,\alpha'(0)\\b\,\beta'(0)
  \end{pmatrix}.
\]
By Lemma~\ref{lem:score-Vi},
$\alpha_i'(0)=p''(\alpha_i)/(2\,p'(\alpha_i))=V_i(\alpha)$,
and similarly for $\beta$ and $\gamma$.
\end{proof}


\subsection{Jacobian contraction}

\begin{proposition}[Jacobian contraction]
\label{prop:jac-contract}
Let $\mathcal{V}:=\{(u,v)\in\R^n\times\R^n:
u^T\mathbf{1}=v^T\mathbf{1}=0\}$.
For all $(u,v)\in\mathcal{V}$:
\begin{equation}\label{eq:contraction}
  \bigl\|J(u,v)\bigr\|^2
  \;\le\;\|u\|^2+\|v\|^2.
\end{equation}
\end{proposition}

\begin{proof}
Fix $(u,v)\in\mathcal{V}$ and consider the paths
$\alpha(t)=\alpha+tu$, $\beta(t)=\beta+tv$,
$\gamma(t)=\Omega(\alpha(t),\beta(t))$.
By variance additivity (Lemma~\ref{lem:var-add}) applied to
root vectors:
\[
  \frac{1}{n}\sum_i\gamma_i(t)^2
  -\Bigl(\frac{1}{n}\sum_i\gamma_i(t)\Bigr)^{\!2}
  =\mathrm{Var}(\alpha(t))+\mathrm{Var}(\beta(t)).
\]
Since $(u,v)\in\mathcal{V}$, the means are constant in~$t$.
Differentiating twice at $t=0$:
\[
  \frac{1}{n}\sum_i\bigl[
    (D_w\gamma_i)^2+\gamma_i\,D_w^2\gamma_i
  \bigr]
  =\frac{1}{n}(\|u\|^2+\|v\|^2),
\]
where $w=(u,v)$, giving the \emph{Hessian identity}
\begin{equation}\label{eq:hessian-id}
  \|Jw\|^2
  =\|u\|^2+\|v\|^2
  -\sum_{i}\gamma_i\,w^T H^{(i)} w,
\end{equation}
with $H^{(i)}:=\mathrm{Hess}\,\Omega_i$ at
$(\alpha,\beta)$.
It remains to show
$\sum_i\gamma_i\,w^T H^{(i)} w\ge 0$.

The polynomial
$f(x,a,b):=\mathbb{E}_{\pi\in S_n}
\prod_{i=1}^n(x-a_i-b_{\pi(i)})$
is hyperbolic in the direction $e_1=(1,0,\ldots,0)$.
By a theorem of Bauschke, G\"uler, Lewis, and
Sendov~\cite{BGLS}, the partial sums
$\sigma_k(a,b):=\sum_{i=1}^k\gamma_i(a,b)$ are convex,
so $\mathrm{Hess}\,\sigma_k\succeq 0$.
Since $\gamma_1\ge\cdots\ge\gamma_n$,
$\sum_i\gamma_i\,H^{(i)}$ is a positive linear combination
of $\mathrm{Hess}\,\sigma_k$, hence
$\sum_i\gamma_i\,w^T H^{(i)} w\ge 0$.
\end{proof}


% ═══════════════════════════════════════════════════════════
\section{Proof of the main theorem}
\label{sec:proof}
% ═══════════════════════════════════════════════════════════

\begin{theorem}\label{thm:full-stam}
Inequality~\eqref{eq:stam} holds for every $n\ge 2$.
\end{theorem}

\begin{proof}
We may assume $p$ and~$q$ have simple roots
(the extension to repeated roots is in
Remark~\ref{rem:boundary}).
Write $\alpha,\beta,\gamma$ for their ordered root vectors.

\emph{Step~1 (Contraction applied to scores).}
By Lemma~\ref{lem:score-ids}(i),
$(V(\alpha),V(\beta))\in\mathcal{V}$.
Apply Proposition~\ref{prop:jac-contract} with
$u=a\,V(\alpha)$, $v=b\,V(\beta)$ ($a,b>0$ to be
chosen):
\[
  \bigl\|J(a\,V(\alpha),\,b\,V(\beta))\bigr\|^2
  \le a^2\,\Phi_n(p)+b^2\,\Phi_n(q).
\]
By Proposition~\ref{prop:score-jac}, the left-hand side
equals $(a+b)^2\,\Phi_n(p\fp q)$.
Hence
\begin{equation}\label{eq:ab-ineq}
  (a+b)^2\,\Phi_n(p\fp q)
  \;\le\;a^2\,\Phi_n(p)+b^2\,\Phi_n(q).
\end{equation}

\emph{Step~2 (Blachman's choice).}
Set $a=1/\Phi_n(p)$ and $b=1/\Phi_n(q)$.
Then~\eqref{eq:ab-ineq} becomes
\[
  \Bigl(\frac{1}{\Phi_n(p)}+\frac{1}{\Phi_n(q)}\Bigr)^{\!2}
  \Phi_n(p\fp q)
  \;\le\;\frac{1}{\Phi_n(p)}+\frac{1}{\Phi_n(q)}.
\]
Dividing both sides by
$\bigl(1/\Phi_n(p)+1/\Phi_n(q)\bigr)\,\Phi_n(p\fp q)>0$
yields~\eqref{eq:stam}.
\end{proof}

\begin{remark}\label{rem:old-proof}
The proof above replaces the case-split argument from
the original version of this paper, in which the forward
bound $1/\Phi_n(p\fp q)\ge(a+b)/(a\,\Phi_n(p))$ and its
reverse were combined.
Those bounds followed from flowing $p$ along~$q$
using fractional $K$-transform powers; however, the
fractional power~$q_t$ is not generally real-rooted
(see Remark~\ref{rem:why-gaussian}),
so the flow does not preserve~$\PnR$ and the root ODE
used in the Lyapunov argument is invalid.
The Jacobian contraction
(Proposition~\ref{prop:jac-contract})
provides the missing bridge.
\end{remark}

\begin{remark}\label{rem:boundary}
With the convention $1/\Phi_n:=0$ for repeated roots,
\eqref{eq:stam} extends to the boundary of~$\PnR$.
When both $p$ and~$q$ have repeated roots, both sides vanish.
When exactly one, say~$p$, has a repeated root, approximate
$p$ by simple-root polynomials $p_\varepsilon\to p$;
the inequality for $p_\varepsilon,q$ gives
$1/\Phi_n(p_\varepsilon\fp q)\ge 1/\Phi_n(q)$,
and continuity of the convolution yields the result in
the limit.
The Nuij operator $T_\varepsilon:=(1-\varepsilon\,d/dx)^n$
provides such approximations explicitly: $T_\varepsilon p$
is monic, real-rooted, and has simple roots for
$\varepsilon>0$; moreover, $T_\varepsilon$ commutes
with~$\fp$, so
$(T_\varepsilon p)\fp(T_\varepsilon q)
=T_\varepsilon^2(p\fp q)$.
\end{remark}


% ═══════════════════════════════════════════════════════════
\section*{References}
% ═══════════════════════════════════════════════════════════

\begin{thebibliography}{9}

\bibitem{BGLS}
H.~H.~Bauschke, O.~G\"uler, A.~S.~Lewis,
and H.~S.~Sendov,
\emph{Hyperbolic polynomials and convex analysis},
Canad.\ J.\ Math.\ \textbf{53} (2001), 470--488.

\bibitem{blachman}
N.~M.~Blachman,
\emph{The convolution inequality for entropy powers},
IEEE Trans.\ Inform.\ Theory \textbf{IT-11} (1965),
267--271.

\bibitem{nuij}
W.~Nuij,
\emph{A note on hyperbolic polynomials},
Math.\ Scand.\ \textbf{23} (1968), 69--72.

\bibitem{MSS}
A.~Marcus, D.~A.~Spielman, and N.~Srivastava,
\emph{Interlacing families {II}: Mixed characteristic
polynomials and the {K}adison--{S}inger problem},
Ann.\ of Math.\ (2) \textbf{182} (2015), 327--350.

\bibitem{stam}
A.~J.~Stam,
\emph{Some inequalities satisfied by the quantities of
information of {F}isher and {S}hannon},
Inform.\ Control \textbf{2} (1959), 101--112.

\end{thebibliography}

\end{document}
